\documentclass[11pt]{letter}
\usepackage{geometry}
\geometry{left=2.5cm, right=2.5cm, top=2.5cm, bottom=2.5cm}
\usepackage{hyperref}
\setlength{\parskip}{0.8em}
\setlength{\parindent}{0pt}
\address{Guangdong University of Finance \\ Guangzhou, China \\ Email: 241734106@m.gduf.edu.cn \\ ORCID: 0009-0008-2164-4557}
\signature{Zhiren Xiao}

\begin{document}

\begin{letter}{Professor Witold Pedrycz \\ Editor-in-Chief \\ \textit{Information Sciences} \\ Elsevier}

\opening{Dear Professor Pedrycz,}

We submit our manuscript entitled ``Exact Bayesian Inference for Spike-and-Slab Priors in Takagi--Sugeno--Kang Fuzzy Systems with Calibrated Model-Averaged Prediction Intervals'' for consideration in \textit{Information Sciences}.

\noindent\textbf{Contribution and Novelty}

This paper presents exact Bayesian inference for spike-and-slab priors in Takagi--Sugeno--Kang (TSK) fuzzy systems, delivering calibrated model-averaged prediction intervals. Whereas prior work has coupled TSK systems either with conjugate priors without sparsity, or with frequentist $L_1$ regularization without uncertainty, we place a spike-and-slab prior on the rule consequents and perform posterior inference by block-Gibbs sampling with Bayesian model averaging (BMA). The method applies to any rule-based regression with linear consequents; TSK fuzzy systems are the concrete instantiation. To our knowledge, no prior work provides exact spike-and-slab posterior sampling with BMA-calibrated prediction intervals for linear-consequent rule models; our method fills this gap. Three results structure the contribution.

\textbf{First}, we provide the exact MCMC machinery: a block-Gibbs sampler whose full conditional for each rule's inclusion indicator is obtained in $\mathcal{O}(n(d+1)^2)$ time via the matrix-determinant and Woodbury identities, together with a BMA predictive distribution whose variance decomposes into aleatoric and model-uncertainty terms by the law of total variance. This yields near-nominal 95\% prediction intervals (PICP $=0.93$--$0.95$) where a BIC-plus-Laplace analytical approximation undercovers (PICP $=0.00$--$0.18$). We trace the failure to the hard-thresholding step and the collapse of the Laplace covariance, providing a constructive diagnosis of when analytical approximations to Bayesian model selection are safe---a question of direct interest to the uncertainty-quantification community.

\textbf{Second}, we document and correct a silent membership-function implementation bug in which fuzzy-set spreads are recomputed from test data at prediction time, corrupting every TSK baseline. Correcting it raises the dense TSK least-squares baseline from $R^2=0.41$ to $0.94$ on the Energy-Cooling target. Because this implementation pattern is one we encountered while reproducing published TSK baselines, we release the corrected code and baselines as a reproducibility benchmark.

\textbf{Third}, we characterize the boundary conditions of sparsity: under correct fuzzy $c$-means, neither rule-level nor coefficient-level sparsity beats the dense baseline on the low-dimensional building-energy domain. We report this as a boundary result, delineating precisely the regime in which sparsity-inducing priors earn their cost.

\noindent\textbf{Journal Fit}

This manuscript's core contribution is a general methodology for calibrated Bayesian inference under model uncertainty---exact block-Gibbs sampling for spike-and-slab priors combined with Bayesian model averaging---with Takagi--Sugeno--Kang fuzzy systems as the concrete instantiation. The topic sits within \textit{Information Sciences}' established treatment of uncertainty: the journal has published Bayesian approaches to uncertainty quantification (e.g., \emph{Combining pre- and post-model information in the uncertainty quantification of non-deterministic models using an extended Bayesian melding approach}, 2019; \emph{Bayesian approach for inconsistent information}, 2013) and counts Zadeh's generalized theory of uncertainty among its canonical references. By contrast, prediction-interval methods for fuzzy systems that remain at the application level have appeared in \emph{Applied Soft Computing}; our contribution is methodological rather than an application of a known technique to building-energy data. We submit it as a contribution to computational intelligence and uncertainty quantification, with the fuzzy-system instantiation as the carrier.

\noindent\textbf{Manuscript Metadata}

\begin{itemize}
\item Figures: 6 (main text)
\item Tables: 1 (main text)
\item References: 38
\item Code and corrected baselines: \url{https://github.com/Jackxiaozhiren/tsk-spikeslab} (persistent DOI: \href{https://doi.org/10.5281/zenodo.21929319}{10.5281/zenodo.21929319})
\end{itemize}

\noindent\textbf{Suggested Reviewers}

\begin{enumerate}
\item Prof.\ Thierry Den\oe ux, Heudiasyc, Universit\'e de Technologie de Compi\`egne, France, \texttt{thierry.denoeux@hds.utc.fr} --- expertise in evidential and uncertainty-aware regression for fuzzy systems.
\item Prof.\ S\'ebastien Destercke, Heudiasyc, CNRS/Universit\'e de Technologie de Compi\`egne, France, \texttt{sebastien.destercke@hds.utc.fr} --- expertise in uncertainty quantification, imprecise probability, and fuzzy systems.
\item Prof.\ Nikhil R. Pal, Electronics and Communication Sciences Unit, Indian Statistical Institute, Kolkata, India, \texttt{nikhil@isical.ac.in} --- expertise in fuzzy systems and high-dimensional TSK fuzzy modeling.
\end{enumerate}

\noindent\textbf{Originality and Ethical Statement}

This manuscript is original, has not been published previously, and is not under consideration elsewhere. The datasets are publicly available from the UCI Machine Learning Repository. The author declares no competing interests.

\closing{Sincerely,}

Zhiren Xiao \\
Guangdong University of Finance, Guangzhou, China

\end{letter}
\end{document}
